\chapter{Uvod}
\label{ch:uvod}

\napomena{Smjernice: Uvod definiše predmet istraživanja i osnovne pojmove, daje teorijski
okvir, ukazuje na nedostatke postojećih rješenja/literature, navodi istraživačka pitanja
i predstavlja strukturu rada. Obim: 3--5 strana. Pisati na kraju, iz nacrta prijave.}

\section{Predmet i cilj rada}
\label{sec:predmet-cilj}
\napomena{Predmet: dvoprolazni asembler za podskup RV32I. Opšti i posebni ciljevi
(preuzeti iz prijave). Ovdje navesti i istraživačka pitanja: jedno glavno i tri pomoćna.}

\section{Motivacija za razvoj asemblerskog prevodioca}
\label{sec:motivacija}
\napomena{Zašto napisati sopstveni asembler: razumijevanje veze hardvera i softvera,
složenost industrijskih alata (GNU as, LLVM), obrazovna vrijednost.}

\section{Značaj asemblera u računarskim sistemima}
\label{sec:znacaj-asemblera}
\napomena{Mjesto asemblera u lancu kompajler → asembler → linker → loader;
primjene: OS, ugrađeni sistemi, kompajleri, reverzni inženjering. \parencite{levine2000}}

\section{Metodologija rada}
\label{sec:metodologija}
\napomena{Metode iz prijave: analiza literature, modelovanje/specifikacija,
konstruktivna metoda, eksperimentalna (testiranje), komparativna, studija slučaja.}

\section{Struktura rada}
\label{sec:struktura}
\napomena{Po jedan pasus za svako poglavlje (\ref{ch:osnove}--\ref{ch:zakljucak}) i priloge.}
